\chapter{Trust, Oracles, and Proofs}

The hybrid system's most distinctive engineering move is to treat semantic judgment as a
controlled effect rather than as an invisible side channel. This is expressed by the
oracle monad
\[
  \TO(A) = A \times \mathsf{TrustLevel} \times \mathsf{Confidence} \times \mathsf{Provenance}.
\]
A semantic check is therefore not merely a boolean returned by an LLM. It is a value with
a trace: what model or process judged it, how confident that judgment is, whether it was
cached, and which trust level should be attached to the result.

\section*{The trust lattice}
The core trust ordering used by the hybrid types can be written as
\[
  \mathrm{CONTRADICTED} < \mathrm{UNTRUSTED} < \mathrm{RUNTIME\_CHECKED} < \mathrm{LLM\_JUDGED} < \mathrm{Z3\_PROVEN} < \mathrm{LEAN\_VERIFIED}.
\]
Some convenience exports use slightly different public names, but the semantic pattern is
the same: trust is ordered, and the order is used algebraically.

The meet operation records weakest-link composition,
\[
  a \sqcap b = \min(a,b),
\]
while the join operation records strongest evidence,
\[
  a \sqcup b = \max(a,b).
\]
These operations matter because they make compositionality honest. If one stage is only
LLM-judged while a downstream stage is solver-proven, the composition should not suddenly
pretend to be Lean-verified. The weakest link constrains the trust of the whole.

\section*{Structural versus semantic predicates}
The oracle monad would be much less useful if every predicate lived inside it. The
repository instead distinguishes structural predicates from semantic ones. Structural
predicates are intended to be decidable or at least deterministically checkable. Semantic
predicates encode claims whose meaningful evaluation may require richer interpretation,
rubrics, or domain knowledge.

This split is the deepest reason the hybrid system can talk about anti-hallucination,
solver discharge, and proofs in one breath without conflating them. Structural checks may
be discharged by Z3 or deterministic evaluation. Semantic checks enter the oracle monad
and produce confidence-bearing evidence. Proof translation may then reflect some of these
facts into Lean, leaving others as assumptions or residual obligations.

\begin{definition}[Hybrid refinement]
A hybrid refinement type is a base type $T$ equipped with a structural predicate
$\varphi_d$ and an optional semantic predicate $\varphi_s$, written
\[
  \{x : T \mid \varphi_d(x) \land \varphi_s(x)\}.
\]
The intent is that $\varphi_d$ is decidable while $\varphi_s$ is judged via the oracle
pipeline and therefore carries trust and confidence metadata.
\end{definition}

\section*{Lean translation and residual honesty}
The repository's Lean translation and discharge cascade are significant because they
preserve the distinction between theorem-backed and oracle-backed content. The discharge
cascade typically moves through trivial discharge, Z3, oracle-assisted proof search,
Lean checking, and residual fallback. This is an unusually honest design. It avoids the
common temptation to treat all successful pipeline stages as equally formal.

A semantic predicate may be accepted at the semantic layer and even exported into a proof
pipeline, but if it ultimately becomes an axiom or a residual assumption that boundary
should remain visible. One of the strongest lessons from the repository is therefore that
proof engineering and trust engineering must coexist. A system that hides its residuals is
less trustworthy than one that openly classifies them.

\begin{proposition}
Any semantics-preserving translation from hybrid types to Lean must either prove,
postulate, or residualize every semantic predicate that lacks a purely structural proof.
\end{proposition}

This proposition is conceptually obvious, but it is exactly the sort of obvious fact that
software tools often ignore in practice. \DepPy{} does not ignore it; the hybrid design is
built around it.
